APR 06 - CODIGO RNN/LSTM/ECG
import numpy as np
from keras.models import Sequential
from keras.layers import Input,LSTM,Dense
from keras.optimizers import Adam
# exemplo de sequencias:
# X.shape = (N, time_steps, features)
net=Sequential([
 Input(shape=(25,1)),
 LSTM(16,return_sequences=True),
 LSTM(8),
 Dense(1,activation="sigmoid")
])
net.compile(optimizer=Adam(learning_rate=.001),
 loss="binary_crossentropy",metrics=["accuracy"])
net.fit(Xtrain,Ttrain,validation_data=(Xval,Tval),
 epochs=50,batch_size=32)
p=(net.predict(Xtest).ravel()>=.5).astype(int)
# ECG: preparar segmentos do mesmo tamanho
FS=300; DUR=30; NS=FS*DUR
sinal=sinal.astype(float)
sinal=(sinal-sinal.mean())/(sinal.std()+1e-8)
# 9000 samples -> (9000,1), depois batch
xraw=sinal[:,None]
# Intervalos RR: indices dos picos R
from scipy.signal import find_peaks,resample_poly
peaks,_=find_peaks(sinal,distance=int(.3*FS))
rr=np.diff(peaks)/FS
hr=60/rr
# reamostrar de 300Hz para 50Hz:
sinal50=resample_poly(sinal,50,300)
# train/test sempre por registo, nao por janela
from sklearn.model_selection import train_test_split
itr,ite=train_test_split(np.arange(len(T)),
 test_size=.3,random_state=42,stratify=T)
# CNN input (N,9000,1)
# LSTM RR input (N,25,1)
# LSTM beats (N,25,181) no exemplo das aulas.